\documentclass{article}
\usepackage{amsmath, amssymb, hyperref}

\title{\flushleft\textbf{Lunch Stealing} \\[1.0em] \small CITS3001 problem \textbf{lunchstealing} - Time limit: 3s \noindent\rule{\textwidth}{0.4pt}}
\date{}

\newcommand{\sampleboxwidth}{\dimexpr\textwidth-2\fboxsep-2\fboxrule\relax}

\begin{document}

\maketitle
\vspace{-90pt}
\section*{}
It is a bright afternoon on the banks of the Derbarl Yerrigan, and the black swans have decided that today is a feast day. Students are sprawled along the river with their lunches, and the flock means to make off with as many as it can.

A swan will snatch a student's lunch --- unless that particular student unnerves it, in which case it keeps a wary distance and leaves that lunch alone. Swans are proud birds: in all the commotion a single swan will carry off at most one lunch, and no lunch is worth squabbling over, so each student's lunch is taken by at most one swan.

For every swan the flock's lookout has noted which students it is too timid to approach. Given those notes, work out the greatest number of lunches the flock can steal between them.

\section*{Input}
The first line contains two integers $S$ and $L$ ($1 \le S, L \le 100$): the number of swans and the number of student lunches. Swans are numbered $1$ to $S$ and lunches $1$ to $L$.

Each of the next $S$ lines describes one swan, in order from swan $1$ to swan $S$. The line begins with the swan's own number, followed by the (distinct) numbers of the students it is too scared to steal from. A line containing only the swan's number means that swan fears no one and could take any lunch.

\section*{Output}
Output a single integer: the maximum number of lunches the flock can steal, given that each swan takes at most one lunch and each lunch is taken by at most one swan.

\section*{Sample Tests}

\subsection*{Sample Input 1}
\fbox{
  \begin{minipage}{\sampleboxwidth}
    \ttfamily
    3 3 \\
    1 3 \\
    2 1 \\
    3 2
  \end{minipage}%
}

\subsection*{Sample Output 1}
\fbox{
  \begin{minipage}{\sampleboxwidth}
    \ttfamily
    3
  \end{minipage}%
}

\noindent Swan 1 fears student 3 (so it may take lunch 1 or 2), swan 2 fears student 1, and swan 3 fears student 2. Assigning swan 1 to lunch 2, swan 2 to lunch 3 and swan 3 to lunch 1 lets all three swans eat.

\subsection*{Sample Input 2}
\fbox{
  \begin{minipage}{\sampleboxwidth}
    \ttfamily
    3 3 \\
    1 3 \\
    2 2 3 \\
    3 1
  \end{minipage}%
}

\subsection*{Sample Output 2}
\fbox{
  \begin{minipage}{\sampleboxwidth}
    \ttfamily
    3
  \end{minipage}%
}

\subsection*{Sample Input 3}
\fbox{
  \begin{minipage}{\sampleboxwidth}
    \ttfamily
    3 3 \\
    1 2 3 \\
    2 2 3 \\
    3 2 3
  \end{minipage}%
}

\subsection*{Sample Output 3}
\fbox{
  \begin{minipage}{\sampleboxwidth}
    \ttfamily
    1
  \end{minipage}%
}

\noindent All three swans are scared of students 2 and 3, so every swan can only approach student 1. Just one of them gets a lunch.

\section*{Note}
A swan that fears every student, or a lunch that no swan dares approach, is possible --- such a swan or lunch simply cannot be part of any steal.

\end{document}
